% AUTORIDAD DE FUENTE:
% expose_VIII_body.tex, líneas 1800--1916; índices de escaneo 289--295;
% folios de fuente 278--284. Las correcciones materiales de la capa
% francesa se conservan reversiblemente mediante \sisi.

\medskip
\noindent\underline{5. Extensiones y biextensiones por esquemas en
grupos constantes truncados sin torsión.}

En esta sección y en la siguiente damos algunos resultados auxiliares
que se utilizarán en la sección 7 en el caso $Z=\mathbb Z$.

\medskip
\noindent\underline{Proposición 5.1.} \underline{Sea $S$ un esquema
irreducible geométricamente unirrama con punto genérico $\eta$, y sea
$Z$ un grupo sin torsión, que define por tanto un esquema en grupos
constante $Z_S$ sobre $S$. Entonces todo torsor $E$ bajo $Z_S$ es
trivial, y el funtor}
\[
 E\longmapsto E(\eta)
\]
\underline{de la categoría de torsors bajo $Z_S$ a la categoría de
torsors bajo $Z=Z_S(\eta)$ es una equivalencia de categorías.}

Un lema conocido de teoría del descenso (SGA 3 X 5.4) muestra que todo
torsor $E$ bajo $\sisi{G_S}{Z_S}$ es representable. De hecho, ese
resultado se aplica a todo esquema en grupos sobre $S$ que sea
separado, localmente de presentación finita y localmente cuasifinito
sobre $S$. Como $Z_S$ es étale sobre $S$, el descenso fpqc muestra que
lo mismo vale para todo torsor $E$ bajo $Z_S$; tal torsor es, por
tanto, localmente trivial para la \underline{topología étale} (y no
solo para la topología fppf).

%%%%%%%%%%  scan idx 290  |  folio 279  %%%%%%%%%%

Para probar que $E$ es trivial basta, en consecuencia, demostrar
\[
 H^1(S_{\mathrm{\acute{e}t}},Z_S)=e.
\]
Del mismo modo, la plena fidelidad de 5.1 se reduce a la relación
\[
 Z\xrightarrow{\ \sim\ }H^0(S_{\mathrm{\acute{e}t}},Z_S).
\]
Esta última es inmediata. Para demostrar la primera, consideremos la
inclusión
\[
 i:\eta\longrightarrow S.
\]
Como $S$ es geométricamente unirrama, tenemos
\[
 Z_S\xrightarrow{\ \sim\ }i_*(Z_\eta)
\]
(SGA 4 IX 2.14.1). La sucesión exacta no conmutativa en grados bajos
que sustituye a la sucesión espectral de Leray para $i:\eta\to S$
(SGA 4 XII 3.2) comienza, por tanto, por
\[
 e\longrightarrow H^1(S,i_*(Z_\eta))
   \longrightarrow H^1(\eta,Z_\eta)\longrightarrow\cdots,
\]
y muestra que
\[
 H^1(S,Z_S)\longrightarrow H^1(\eta,Z_\eta)
\]
es inyectiva. Quedamos así reducidos al caso en que $S$ es el espectro
de un cuerpo $k$. Si $\pi$ es el grupo de Galois de una clausura
separable de $k$, entonces
\[
 H^1(k,Z)=\Hom_{\mathrm{cont}}(\pi,Z)/Z=e,
\]
porque $\sisi{E}{Z}$ no tiene torsión. Esto concluye la demostración.

\medskip
\noindent\underline{Corolario 5.2.} \underline{Con la notación de
5.1, sea $P$ un esquema liso sobre $S$. Entonces todo torsor $E$ bajo
$Z_P$ es trivial, y el funtor de restricción}
\[
 E\longmapsto E_\eta
\]
\underline{de la categoría de torsors sobre
$\sisi{P_\eta}{P}$ bajo $\sisi{Z_{P_\eta}}{Z_P}$ a la categoría de
torsors sobre $P_\eta$ bajo $Z_{P_\eta}$ es una equivalencia de
categorías.}

%%%%%%%%%%  scan idx 291  |  folio 280  %%%%%%%%%%

En efecto, $P$ es geométricamente unirrama (EGA IV 17.5.7). Además,
cada componente irreducible de $P$ domina $S$, porque $P$ es plano
sobre $S$, y el conjunto de componentes irreducibles de $P$ es
localmente finito. Así, $P$ es la suma disjunta de esquemas irreducibles
geométricamente unirrama $P_i$. La afirmación se deduce inmediatamente
de 5.1 aplicado a los distintos $P_i$ y a sus fibras genéricas
$P_{i\eta}$, observando que el punto genérico $p_i$ de $P_i$ es también
el punto genérico de $P_{\sisi{i}{i\eta}}$.

\medskip
\noindent 5.2.1. El corolario precedente da inmediatamente la
afirmación análoga para bitorsors sobre $P$ bajo $(Z_P,Z_P)$: la
restricción a la fibra genérica induce una equivalencia de su categoría
con la categoría de bitorsors sobre $P_\eta$ bajo
$(Z_{P_\eta},Z_{P_\eta})$. En efecto, basta expresar la estructura de
bitorsor sobre $E$ mediante su estructura de torsor derecho bajo $Z_P$
y un isomorfismo de $Z_P$ con el conmutante de $Z_P$ para su acción
sobre $E$. Puesto que $E$ es necesariamente trivial, dicho conmutante
es isomorfo a $Z_P$. Quedamos reducidos, en cada componente irreducible
$X$ de $P$, a la observación inmediata siguiente: si $X$ es un esquema
conexo no vacío, el funtor $Z\mapsto Z_X$ de conjuntos a esquemas sobre
$X$ es plenamente fiel.

%%%%%%%%%%  scan idx 292  |  folio 281  %%%%%%%%%%

\medskip
\noindent\underline{Proposición 5.3.} \underline{Sean $S$ y $Z$ como
en 5.1, y sea $P$ un esquema en grupos liso sobre $S$. Consideremos
extensiones $E$ de esquemas en grupos de $P$ por $Z_S$. Entonces la
restricción a la fibra genérica}
\[
 E\longmapsto E_\eta
\]
\underline{es una equivalencia de la categoría de estas extensiones
con la categoría de extensiones de $P_\eta$ por $Z_\eta$. Si $P$ y $Z$
son conmutativos, induce también una equivalencia entre las categorías
correspondientes de extensiones conmutativas.}

Observemos primero que toda extensión de haces de grupos de $P$ por
$Z_S$ (y, análogamente, de $P_\eta$ por $Z_\eta$) da, por 5.2, un torsor
trivial y, en particular, representable. Las categorías de 5.3 pueden,
por tanto, interpretarse igualmente como categorías de extensiones de
haces de grupos. La descripción de VII 1.1.6 y 1.2 las expresa mediante
bitorsors sobre $P$, $P\times P$ y $P\times P\times P$
(respectivamente sobre $P_\eta$, $P_\eta\times P_\eta$ y
$P_\eta\times P_\eta\times P_\eta$), bajo las parejas de grupos
obtenidas de $(Z,Z)$ por cambio de base. Aplicando 5.2 en el caso
conmutativo y 5.2.1 en el caso general a los esquemas lisos sobre $S$
$P$, $P\times P$ y $P\times P\times P$, se obtiene la conclusión
afirmada. El mismo argumento da el resultado siguiente.

\medskip
\noindent\underline{Corolario 5.4.} \underline{Sean $S,Z$ como en
5.1, y sean $P,Q$ dos esquemas en grupos lisos sobre $S$. Supongamos
que $P,Q$ y $Z$ son conmutativos, y consideremos biextensiones $E$ de
$(P,Q)$ por $Z_S$. Entonces la restricción a la fibra genérica}
\[
 E\longmapsto E_\eta
\]
\underline{es una equivalencia de la categoría de estas biextensiones
con la categoría de biextensiones de $(P_\eta,Q_\eta)$ por $Z_\eta$.}

%%%%%%%%%%  scan idx 293  |  folio 282  %%%%%%%%%%

Los resultados precedentes reducen la determinación de las categorías
de extensiones o biextensiones por un grupo constante sin torsión al
caso en que el esquema base es el espectro de un cuerpo. En ese caso
obtenemos lo siguiente.

\medskip
\noindent\underline{Proposición 5.5.} \underline{Sea $k$ un cuerpo,
$Z$ un grupo sin torsión, $P$ un esquema en grupos sobre $k$ localmente
de tipo finito, $P^\circ$ su componente neutra y
$\Phi=P/P^\circ$ su grupo de componentes conexas (SGA 3
$\mathrm{VI}_A$ 4). Entonces:}

\smallskip
\noindent (i) \underline{El funtor natural de la categoría de esquemas
en grupos que son extensiones (respectivamente, extensiones
conmutativas) de $\Phi$ por $Z_k$ a la categoría de esquemas en grupos
que son extensiones (respectivamente, extensiones conmutativas) de $P$
por $Z_k$ es una equivalencia. En particular, si $P$ es conexo, toda
extensión de $P$ por $Z_k$ es trivial y la categoría de tales
extensiones es rígida.}

\smallskip
\noindent (ii) \underline{Supongamos que $\Phi$ es un grupo de torsión
y que $P$ y $Z$ son conmutativos. Entonces la categoría de extensiones
de $P$ por $Z_k$ es rígida: todo automorfismo de un objeto de esta
categoría es trivial. Si además existe un entero $n>0$ tal que
$n\id_\Phi=0$ —por ejemplo, si $\Phi$ es finito, como ocurre cuando
$\sisi{Z}{P}$ es de tipo finito—, entonces el grupo de clases de
extensiones conmutativas de $P$ por $\sisi{G_k}{Z_k}$ viene dado por el
isomorfismo canónico}
\begin{equation}\tag{5.5.1}
 \Ext^1(P,Z_k)
 \xleftarrow{\ \sim\ }
 \Hom\bigl(\Phi,(Z\otimes_{\mathbb Z}\mathbb Q)/Z\bigr),
\end{equation}
\underline{definido por el operador de borde de la sucesión exacta de
grupos $\Ext$ asociada a la sucesión exacta}
\begin{equation}\tag{5.5.2}
 0\longrightarrow Z\longrightarrow Z\otimes_{\mathbb Z}\mathbb Q
 \longrightarrow (Z\otimes_{\mathbb Z}\mathbb Q)/Z
 \longrightarrow0.
\end{equation}

%%%%%%%%%%  scan idx 294  |  folio 283  %%%%%%%%%%

El morfismo canónico
\[
 P\longrightarrow\Phi
\]
induce manifiestamente un isomorfismo entre los conjuntos de
componentes irreducibles, pues es plano con fibras irreducibles. Por
5.1, el funtor inducido de torsors bajo $Z_\Phi$ a torsors bajo $Z_P$
es, por tanto, una equivalencia. Recordemos que un esquema en grupos
localmente de tipo finito sobre $k$ es geométricamente unirrama. La
misma conclusión vale para los bitorsors y para los morfismos
\[
 P\times P\longrightarrow\Phi\times\Phi,
 \qquad
 P\times P\times P\longrightarrow\Phi\times\Phi\times\Phi.
\]
La descripción de VII 1.1.6 de las extensiones de $P$ por $Z_k$ y de
$\Phi$ por $Z_k$ (y VII 1.2 en el caso conmutativo) da las afirmaciones
de (i).

Restrinjámonos a las extensiones conmutativas. El grupo de
automorfismos de una extensión de $P$ por $\sisi{G_k}{Z_k}$ es
\[
 \Hom(P,\sisi{G_k}{Z_k})\simeq
 \Hom(\Phi,\sisi{G_k}{Z_k}),
\]
que es manifiestamente nulo si $\Phi$ es un grupo de torsión. Del mismo
modo, si $V=\sisi{G}{Z}\otimes_{\mathbb Z}\mathbb Q$ y $\Phi$ es
anulado por el entero $n>0$, entonces
\[
 \Ext^i(\Phi,V_k)=0\qquad\text{para todo }i\in\mathbb Z,
\]
porque los grupos en cuestión son espacios vectoriales sobre
$\mathbb Q$ anulados por $n$. El isomorfismo (5.5.1) se sigue
inmediatamente.

El mismo método demuestra, esencialmente, lo siguiente.

\medskip
\noindent\underline{Corolario 5.6.} \underline{Sea $k$ un cuerpo, sea
$Z$ un grupo conmutativo sin torsión, sean $P$ y $Q$ dos esquemas en
grupos conmutativos localmente de tipo finito sobre $k$, y sean
$\Phi=P/P^\circ$ y $\Psi=Q/Q^\circ$ sus grupos de componentes
$\sisi{neutras}{conexas}$. Entonces:}

%%%%%%%%%%  scan idx 295  |  folio 284  %%%%%%%%%%

\smallskip
\noindent (i) \underline{El funtor natural de la categoría de
biextensiones de $(\Phi,\Psi)$ por $Z_k$ a la categoría de
biextensiones de $(P,Q)$ por $Z_k$ es una equivalencia. En particular,
si $P$ o $Q$ es conexo, la categoría de biextensiones de $(P,Q)$ por
$Z_k$ es equivalente a la categoría terminal.}

\smallskip
\noindent (ii) \underline{Supongamos que $\Phi$ o $\Psi$ es un grupo
de torsión. Entonces la categoría de biextensiones de $(P,Q)$ por
$Z_k$ es rígida. Si además existe un entero $n>0$ tal que
$n\id_\Phi=0$ o $n\id_\Psi=0$, entonces el grupo de clases de
biextensiones de $(P,Q)$ por $Z_k$ viene dado por el isomorfismo
canónico}
\begin{equation}\tag{5.6.1}
 \Biext^1(P,Q;\sisi{G_k}{Z_k})
 \xleftarrow{\ \sim\ }
 \Hom\bigl(\Phi\otimes\Psi,
 (Z\otimes_{\mathbb Z}\mathbb Q)/Z\bigr),
\end{equation}
\underline{definido por el operador de borde de la sucesión exacta de
grupos $\Biext^i$ $(i=0,1)$ asociada a (5.5.2).}
